% Generated draft: requires mathematical and source review.
\documentclass[11pt]{article}
\usepackage[margin=25mm]{geometry}
\usepackage{fontspec}
\setmainfont{Latin Modern Roman}
\usepackage{amsmath,amssymb,mathtools,unicode-math}
\setmathfont{Latin Modern Math}
\usepackage{enumitem,needspace,xurl,hyperref}
\setlength{\parindent}{0pt}
\setlength{\parskip}{5pt}
\setlength{\emergencystretch}{3em}
\newcommand{\N}{\mathbb{N}}
\newcommand{\Z}{\mathbb{Z}}
\newcommand{\Q}{\mathbb{Q}}
\newcommand{\R}{\mathbb{R}}
\newcommand{\C}{\mathbb{C}}
\newcommand{\F}{\mathbb{F}}
\newcommand{\A}{\mathbb{A}}
\newcommand{\PP}{\mathbb{P}}
\newcommand{\E}{\mathbb{E}}
\newcommand{\eps}{\varepsilon}
\newcommand{\dd}{\,\mathrm{d}}
\begin{document}
\section*{UnsolvedMath ID: 11000020}
\textbf{Source problem code:} catalogue entry
\subsection{Definitions and mathematical statement}
Let \(\mathcal M_g\) be the moduli space of compact genus-\(g\) Riemann surfaces up to biholomorphism, \(g\ge2\), and let \(\operatorname{Aut}(X)\) be the finite holomorphic automorphism group. A group is nilpotent if its lower central series terminates at the trivial group. Problem 2.19 seeks properties of automorphisms or automorphism groups that single out a unique point of \(\mathcal M_g\), providing canonical basepoints in that sense.

Its specific example asks whether there is a unique surface whose full automorphism group is nilpotent of maximum possible order among such surfaces. Set
\[
 \mathcal N_g=\{[X]\in\mathcal M_g:\operatorname{Aut}(X)\text{ is nilpotent}\},
 \qquad m_g=\max_{[X]\in\mathcal N_g}|\operatorname{Aut}(X)|,
\]
when this maximum is defined. The uniqueness question is
\[
 \left|\{[X]\in\mathcal N_g:|\operatorname{Aut}(X)|=m_g\}\right|=1\ ?
\]
The maximum is intended genus by genus. The source does not ask merely for a largest nilpotent subgroup inside an arbitrary, possibly nonnilpotent, full automorphism group; that alternative interpretation would change the extremal set and requires source confirmation if intended.

The surrounding text defines a Hurwitz surface by
\[
 |\operatorname{Aut}(X)|=84(g-1).
\]
It states that the quotient is the \((2,3,7)\) orbifold, with unique hyperbolic metric, and that for each fixed \(g\) there are finitely many such surfaces. It then asks for their exact number
\[
 h_g=\left|\{[X]\in\mathcal M_g:|\operatorname{Aut}(X)|=84(g-1)\}\right|.
\]
No existence in every genus is asserted; \(h_g\) may be zero. The broad canonical-property search, the nilpotent-maximizer uniqueness question and the Hurwitz counting question are retained as distinct tasks. The source supplies no proposed list of canonical properties or formula for either count.
\subsection{Short English statement}
Find automorphism-theoretic conditions selecting unique genus-\(g\) surfaces, especially a unique maximizer with nilpotent full automorphism group. Also count Hurwitz surfaces in each genus. These are uniqueness and counting questions in moduli, not an assertion that a maximally symmetric surface always exists or is unique.
\subsection{Sources}
\begin{itemize}
\item Local UnsolvedMath registry, permanent ID \(11000020\), exact source code \texttt{catalogue entry}, Problem 2.19, ``Canonical basepoints for \(\mathcal M_g\),'' PDF page 30. No individual proposer or proposed year is supplied.
\item Farb, edited volume \emph{Problems on Mapping Class Groups and Related Topics}: \url{https://www.math.uchicago.edu/~farb/papers/mcgbook.pdf}. The volume has 329 explicitly typeset Problem/Conjecture/Question headings, with repeated local numbers distinguished by PDF page.
\item Extraction pdf-font-structure, accessed 2026-07-29. Source evidence NEEDS\_REVIEW presents an open problem and calls for a current primary-literature audit. Supplied OPEN-TRIAGE summary says the canonical-automorphism and maximal-nilpotent uniqueness questions appear open in general; it supplies no separate references or complete Hurwitz-count status assessment.
\item Dataset attribution: Ulam AI / UnsolvedMath Contributors, CC BY 4.0. Source rights NEEDS\_REVIEW: publicly accessible author/publisher copy, redistribution terms requiring review.
\end{itemize}
\end{document}
